% Authoritative LaTeX; edit this file, not the historical Markdown.
\chapter{Tensors Without Magic}\label{chapter-5-tensors-without-magic}
Part I ended with a complete loop. Text became token IDs. A tiny
language model turned the last ID into an embedding, multiplied that
activation by an output projection, produced logits, and fed a selected
token back into the next step. Every number was small enough to print.
Every array was short enough to count.

That convenience hid a systems problem. The model stored its embedding
in a \texttt{Vec\textless{}f32\textgreater{}}, yet the vector itself did
not say that its logical shape was \texttt{{[}vocabulary,\ hidden{]}}.
The projection used another \texttt{Vec\textless{}f32\textgreater{}}. A
comment and a pair of loop bounds carried the missing meaning. If those
facts ever disagreed, the bytes could remain finite and the program
could still compute the wrong answer.

This chapter gives those bytes a contract. It does not build a tensor
framework or perform matrix multiplication. It builds the smallest
substrate on which the rest of Part II can reason honestly about memory.

\begin{quote}
\textbf{FIRST PRINCIPLE} Mathematics tells us \emph{what} value we need.
Tensor layout tells the machine \emph{where} that value lives.
\end{quote}

\section{Where we are: from a model equation to six stored
values}\label{where-we-are-from-a-model-equation-to-six-stored-values}

Follow the engine downward: a request invokes a model; the model asks an
operator to read tensors; the operator must find the right bytes. This
chapter works at that final representation boundary. Chapter 6 will put
arithmetic on top of it. The current repository already contains later
chapters, but this chapter's executable contract remains Tensor
Substrate v1.

\EngineFigure{tensor}{Logical coordinates reach physical storage through explicit strides and a base offset.}{fig:tensor}

Our recurring matrix is \(\mathbf{A}\in\mathbb{R}^{2\times3}\), with
first row \texttt{{[}1,2,3{]}} and second row \texttt{{[}4,5,6{]}}. Its
owner holds six F32 elements, or 24 payload bytes. The first plate
connects the logical coordinate \texttt{{[}1,2{]}} to element offset 5
and byte displacement 20. Neither number is an allocator address. The
actual allocation may live anywhere suitable for \texttt{f32}.

Watch three things separately as the chapter proceeds: which logical
value a coordinate selects, which storage slot supplies it, and which
object owns that slot. Transpose changes the first relationship without
moving the payload. Materialization creates a second owner. Reshape can
have the same output shape as transpose while selecting different
values. Those differences are the mechanism, not details to hide under a
framework method name.

The plates have
\href{https://github.com/hermonai/inside-the-llm-engine/blob/astra-visual-rewrite/figures/chapter05-atlas.md}{plain-text
equivalents} and a
\href{https://github.com/hermonai/inside-the-llm-engine/blob/astra-visual-rewrite/code/reference/fixtures/chapter05-visual.json}{shared
executable fixture}. The original letter-valued and higher-rank
exercises remain useful independent checks. The main visual journey now
uses the same six numbers throughout.

\section{The vector that could not tell the
truth}\label{the-vector-that-could-not-tell-the-truth}

Suppose a loader hands us twelve \texttt{f32} values. What tensor did it
load?

\begin{verbatim}
[12]       [3, 4]       [4, 3]       [2, 2, 3]       [1, 3, 4]
\end{verbatim}

Every listed shape has twelve elements. They do not have the same axes
or indexing semantics. Even shape \texttt{{[}3,4{]}} does not tell us
whether adjacent bytes move across a row, down a column, or through some
strided view of a larger allocation. Twelve scalar payloads are evidence
about storage length, not a complete tensor description.

For every tensor in this book, ask five questions:

\begin{enumerate}
\def\labelenumi{\arabic{enumi}.}
\tightlist
\item
  What is its shape?
\item
  What is its dtype?
\item
  What is its physical layout?
\item
  Who owns the underlying bytes?
\item
  How long may a view reference those bytes?
\end{enumerate}

Later we will add location---CPU memory, device memory, mapped file, or
another tier. Chapter 5 stays in local CPU memory so we can settle the
first five without hiding them behind a device abstraction.

\section{From scalar to tensor}\label{from-scalar-to-tensor}

A \textbf{scalar} is one value and has rank 0. A \textbf{vector} has one
logical axis and rank 1. A \textbf{matrix} has two axes and rank 2. We
use \textbf{tensor} for the general case: a typed multidimensional array
with explicit shape and layout semantics.

{\def\LTcaptype{none} % do not increment counter
\begin{longtable}[]{@{}
  >{\raggedright\arraybackslash}p{(\linewidth - 6\tabcolsep) * \real{0.2308}}
  >{\raggedright\arraybackslash}p{(\linewidth - 6\tabcolsep) * \real{0.2308}}
  >{\raggedright\arraybackslash}p{(\linewidth - 6\tabcolsep) * \real{0.2308}}
  >{\raggedleft\arraybackslash}p{(\linewidth - 6\tabcolsep) * \real{0.3077}}@{}}
\toprule\noalign{}
\begin{minipage}[b]{\linewidth}\raggedright
Object
\end{minipage} & \begin{minipage}[b]{\linewidth}\raggedright
Example
\end{minipage} & \begin{minipage}[b]{\linewidth}\raggedright
Shape
\end{minipage} & \begin{minipage}[b]{\linewidth}\raggedleft
Tensor rank
\end{minipage} \\
\midrule\noalign{}
\endhead
\bottomrule\noalign{}
\endlastfoot
Scalar & \texttt{42.0} & \texttt{{[}{]}} & 0 \\
Vector & \texttt{{[}1,2,3{]}} & \texttt{{[}3{]}} & 1 \\
Matrix & three rows, two columns & \texttt{{[}3,2{]}} & 2 \\
Higher-rank tensor & two batches of three-by-four values &
\texttt{{[}2,3,4{]}} & 3 \\
\end{longtable}
}

Here \textbf{rank} means number of logical axes. Linear algebra also
uses \emph{matrix rank} for the dimension of a matrix's row or column
space. These are different concepts. A shape \texttt{{[}3,2{]}} tensor
always has tensor rank 2, even if all its rows are linearly dependent.

If a rank-\(r\) tensor has shape

\begin{equation}
\mathbf{d} = [d_0,d_1,\ldots,d_{r-1}],
\end{equation}

then \(d_a\) is the length of axis \(a\). A later activation might use
named axes such as \texttt{{[}batch,\ token,\ hidden{]}}; the names
explain model meaning while the integer dimensions define valid indices.

\section{Element count is checked
arithmetic}\label{element-count-is-checked-arithmetic}

The number of logical elements is the product of all dimensions:

\begin{equation}
N(\mathbf{d}) = \prod_{a=0}^{r-1} d_a.
\end{equation}

For shape \texttt{{[}3,4{]}}, \(N=12\). For rank 0, the empty product is
one, so shape \texttt{{[}{]}} describes one scalar. This definition is
mathematically compact, but a systems implementation must ask whether
the product fits its integer type.

\begin{verbatim}
shape.iter().try_fold(1_usize, |count, &dimension| {
    count.checked_mul(dimension).ok_or(TensorError::ShapeOverflow)
})
\end{verbatim}

The real implementation handles zero dimensions first, then uses checked
multiplication. It also checks the byte count:

\begin{equation}
B=N(\mathbf{d})\,w_{\mathrm{dtype}}\quad[\mathrm{bytes}],
\end{equation}

where \(w_{\mathrm{dtype}}\) is bytes per stored element.

Here \(B\) is bytes, not elements. For \texttt{{[}2,3,4{]}} in
\texttt{f32}, \(N=24\) elements and \(B=24\times4=96\) bytes. Both
products can overflow independently.

\begin{quote}
\textbf{ENGINEERING FAILURE --- WRAPPED SIZE} External metadata claims
shape \texttt{{[}usize::MAX,2{]}}. Unchecked multiplication wraps to a
small value, code allocates too little storage, and later address
arithmetic trusts the original dimensions. Checked size arithmetic is
not defensive decoration; it is part of model-loading security.
\end{quote}

\section{Zero-sized dimensions are
deliberate}\label{zero-sized-dimensions-are-deliberate}

Tensor Substrate v1 permits zero-sized dimensions. Shape
\texttt{{[}0,4{]}} has zero elements and owns an empty vector. It has
rank 2, canonical strides \texttt{{[}4,1{]}}, and no valid logical index
because index 0 is already outside axis 0.

This decision keeps shape algebra honest. Empty batches and empty slices
can be represented without a sentinel. A zero-element view may have its
base offset at the end of storage, because it reaches no element. It may
reshape only to a shape whose checked element count is also zero.

Allowing empty tensors does not disable all validation. Canonical stride
products still have to fit \texttt{usize}; malformed metadata cannot
smuggle an unrepresentable layout through a zero dimension.

\section{Dtype is interpretation, not
width}\label{dtype-is-interpretation-not-width}

A \textbf{dtype} specifies how stored bits represent values and
participate in arithmetic. Tensor Substrate v1 implements only
\texttt{F32}. Each element is one IEEE-754 binary32 value and occupies
four bytes, while reductions in later operators may choose a different
accumulation type.

Other systems use \texttt{f16}, \texttt{bf16}, integers, and packed
quantized formats. Two types can occupy the same byte width yet
interpret bits differently. A packed quantized block may represent many
logical weights with shared scale metadata, so ``one scalar equals one
fixed byte range'' eventually stops being a useful model. Chapter 16
will face that complexity. Adding a generic dtype hierarchy now would
create machinery without a chapter requirement.

\texttt{DType::F32} is still explicit in v1. That apparently redundant
name prevents future code from treating
\texttt{\seqsplit{size\_of::<f32>()}} as the whole semantic contract.

\section{Logical values and physical
storage}\label{logical-values-and-physical-storage}

Consider this logical matrix:

\begin{equation}
\mathbf{A} =
\begin{bmatrix}
1 & 2 & 3 \\
4 & 5 & 6
\end{bmatrix},
\qquad \mathbf{A}\in\mathbb{R}^{2\times3}.
\end{equation}

A canonical row-major representation places each row consecutively:

\begin{equation}\mathrm{storage}(A)=(1,2,3,4,5,6),\qquad \mathrm{offset}(i,j)=3i+j.\end{equation}

The complete canonical diagram is
\href{https://github.com/hermonai/inside-the-llm-engine/blob/astra-visual-rewrite/diagrams/tensor/logical-vs-physical.txt}{logical
tensor versus physical storage}. The logical coordinate
\texttt{{[}1,1{]}} names value 5; the layout maps it to physical offset
4. \textbf{Shape} determines whether \texttt{{[}1,1{]}} is a legal
coordinate. \textbf{Strides} determine how that coordinate moves through
storage.

Column-major storage is another valid convention. It would place
\texttt{1,4,2,5,3,6} consecutively for the same logical matrix.
Row-major is v1's chosen owner layout, not a claim that one convention
is universally faster.

\section{Strides turn indices into
offsets}\label{strides-turn-indices-into-offsets}
\index{tensor!strides}

A stride is the distance in storage caused by incrementing one logical
axis by one. Tensor Substrate v1 measures strides in \textbf{elements}.
NumPy and GGML expose byte strides, so source code crossing those
boundaries must convert units explicitly.

Let a view have shape \(\mathbf d\), strides \(\mathbf s\), base element
offset \(b\), and logical index \(\mathbf i\). Its physical element
offset is

\begin{equation}
\operatorname{offset}(\mathbf i)
= b + \sum_{a=0}^{r-1} i_a s_a,
\qquad 0 \le i_a < d_a.
\end{equation}

Every symbol has a physical interpretation:

\begin{itemize}
\tightlist
\item
  \(b\) is an element index into the borrowed slice;
\item
  \(i_a\) is a logical coordinate with no byte unit;
\item
  \(s_a\) is elements per one-axis step;
\item
  the result is an element offset, not a pointer and not a byte address.
\end{itemize}

For contiguous row-major shape \texttt{{[}2,3,4{]}}, calculate strides
from right to left:

\begin{equation}
s_{r-1}=1,
\qquad
s_a=\prod_{k=a+1}^{r-1}d_k.
\end{equation}

The result is \texttt{{[}12,4,1{]}}. Index \texttt{{[}1,2,3{]}} maps to

\begin{equation}
1\times12 + 2\times4 + 3\times1 = 23\text{ elements}.
\end{equation}

Because the dtype is \texttt{f32}, the illustrative byte displacement is
\(23\times4=92\) bytes. See the full
\href{https://github.com/hermonai/inside-the-llm-engine/blob/astra-visual-rewrite/diagrams/tensor/row-major-offsets.txt}{offset
derivation}.

\section{One authoritative checked
offset}\label{one-authoritative-checked-offset}

V1 sends every checked read through \texttt{checked\_offset}. The
operation follows a fixed order:

\begin{enumerate}
\def\labelenumi{\arabic{enumi}.}
\tightlist
\item
  Validate that shape and stride ranks match.
\item
  Validate that the metadata's reachable extent fits storage.
\item
  Validate that index rank equals tensor rank.
\item
  Validate each index against its dimension.
\item
  Multiply index by stride with \texttt{checked\_mul}.
\item
  Add each contribution and the base with \texttt{checked\_add}.
\item
  Perform the final safe slice lookup.
\end{enumerate}

Validating indices before accepting their contribution matters.
Computing an offset first and checking only the final number can let a
malformed coordinate alias an apparently valid physical location.
Logical bounds and physical bounds are separate invariants.

The public API returns
\texttt{Result\textless{}\&f32,\ TensorError\textgreater{}} rather than
implementing Rust's \texttt{Index} trait. \texttt{Index} conventionally
panics on failure; a typed result makes model-derived metadata and
indices inspectable without a process-level surprise.
\texttt{\seqsplit{get2(row,column)}} is a convenience that calls the
same general path.

\section{Storage extent is not element
count}\label{storage-extent-is-not-element-count}

An arbitrary strided view can touch more storage than its logical
element count suggests. Shape \texttt{{[}2,2{]}} has four logical
values. With strides \texttt{{[}100,1{]}}, the view reaches offsets 0,
1, 100, and 101. Four backing elements are nowhere near enough.

For nonempty v1 views with non-negative strides, the largest reachable
offset and required storage length are

\begin{equation}
\begin{aligned}
o_{\max}&=b+\sum_{a=0}^{r-1}(d_a-1)s_a,\\
L_{\min}&=o_{\max}+1
=1+b+\sum_{a=0}^{r-1}(d_a-1)s_a
\quad[\mathrm{elements}].
\end{aligned}
\end{equation}

Every multiply and add is checked. The constructor requires
\(L_{\min}\le L_{\text{storage}}\). For a zero-element view there is no
reachable offset, so the rule becomes \(b\le L_{\text{storage}}\). The
\href{https://github.com/hermonai/inside-the-llm-engine/blob/astra-visual-rewrite/diagrams/tensor/storage-extent.txt}{storage-extent
diagram} shows the smallest counterexample to validating a view by
element count alone.

This formula relies on non-negative strides. V1 uses \texttt{usize} and
does not support reverse traversal through negative strides. It accepts
zero strides for immutable views: several logical indices may then read
the same physical value. That is safe for reading but not canonical
contiguous storage.

\section{Contiguous means one strict thing
here}\label{contiguous-means-one-strict-thing-here}
\index{tensor!contiguity}

A tensor is often called contiguous when logical iteration visits a
dense region in an expected order. Mature frameworks account for memory
formats and size-one dimensions in more flexible ways. Tensor Substrate
v1 intentionally uses a narrower definition:

\begin{quote}
A view is canonical row-major contiguous exactly when its stride vector
equals \texttt{\seqsplit{canonical\_row\_major\_strides(shape)}}.
\end{quote}

Thus \texttt{{[}2,3{]}\ /\ {[}3,1{]}} is contiguous. Its transpose
\texttt{{[}3,2{]}\ /\ {[}1,3{]}} is not. Shape \texttt{{[}1,3,1{]}} has
canonical strides \texttt{{[}3,1,1{]}}; a different stride on a
length-one axis is non-canonical even if it reaches the same values.
This strict rule gives Chapter 6 a simple kernel precondition. We
describe it as \emph{v1 canonical contiguity}, not universal framework
semantics.

The
\href{https://github.com/hermonai/inside-the-llm-engine/blob/astra-visual-rewrite/diagrams/tensor/contiguous-vs-strided.txt}{contiguous
versus strided} diagram shows why a valid logical walk need not be
physical-order traversal.

\section{A view owns metadata, not
elements}\label{a-view-owns-metadata-not-elements}

\texttt{OwnedTensor} owns \texttt{Vec\textless{}f32\textgreater{}},
shape, and validated canonical strides.
\texttt{TensorView\textless{}\textquotesingle{}a\textgreater{}} owns
shape/stride/base metadata but borrows
\texttt{\&\textquotesingle{}a\ {[}f32{]}}. Creating a view may allocate
those small metadata vectors; it never allocates or copies the element
payload.

\begin{verbatim}
pub struct TensorView<'a> {
    storage: &'a [f32],
    shape: Vec<usize>,
    strides: Vec<usize>,
    base_offset: usize,
}
\end{verbatim}

The lifetime \texttt{\textquotesingle{}a} ties the view to valid
storage. It cannot outlive its owner. The compiler rejects returning a
view after a locally created owner has been dropped. This is not merely
a language lesson: a dangling model-weight view is a use-after-free at
the heart of numerical execution.

Multiple immutable views may overlap. A source view, transpose, and row
slice can all read some of the same values. They \textbf{alias} because
their reachable storage regions overlap. The
\href{https://github.com/hermonai/inside-the-llm-engine/blob/astra-visual-rewrite/diagrams/tensor/tensor-ownership.txt}{ownership}
and
\href{https://github.com/hermonai/inside-the-llm-engine/blob/astra-visual-rewrite/diagrams/tensor/tensor-memory-lifetime.txt}{lifetime}
diagrams make those relationships explicit.

\EngineFigure{tensor-ownership}{Owned fields and borrowed slices have different relationships in the actual Rust structs.}{fig:tensor-ownership}

The filled diamond is UML composition: \texttt{OwnedTensor} contains a
\texttt{Vec\textless{}f32\textgreater{}} that owns initialized element
storage. The dashed dependencies are borrows of that storage, not
inheritance. Each view also owns its shape and stride vectors;
\texttt{TensorView} additionally stores a base offset.
\texttt{TensorViewMut} has no arbitrary-stride constructor and no
base-offset field because v1 limits it to the complete canonical slice.
The drawing shows the two borrowing alternatives, not permission to use
an overlapping shared view during exclusive access.

Rust lifetimes describe permitted use, not an obligation to keep a
borrow active until the end of a lexical block. The last use of a shared
view can precede a mutable borrow even when its variable remains in
scope. A surviving reference to one of its elements still counts as a
use of the same borrow. Ownership is therefore about the whole reference
chain, not merely the name of the view object.
\href{https://github.com/hermonai/inside-the-llm-engine/blob/astra-visual-rewrite/figures/generated/tensor-ownership.txt}{Text/UML
source}.

\section{Copy means a new owner}\label{copy-means-a-new-owner}

A copy creates new element storage. It has a separate lifetime, separate
mutation behavior, and a memory-traffic cost proportional to copied
payload. V1 names that boundary \texttt{to\_contiguous}.

{\def\LTcaptype{none} % do not increment counter
\begin{longtable}[]{@{}
  >{\raggedright\arraybackslash}p{(\linewidth - 4\tabcolsep) * \real{0.3333}}
  >{\raggedright\arraybackslash}p{(\linewidth - 4\tabcolsep) * \real{0.3333}}
  >{\raggedright\arraybackslash}p{(\linewidth - 4\tabcolsep) * \real{0.3333}}@{}}
\toprule\noalign{}
\begin{minipage}[b]{\linewidth}\raggedright
Operation
\end{minipage} & \begin{minipage}[b]{\linewidth}\raggedright
Element allocation
\end{minipage} & \begin{minipage}[b]{\linewidth}\raggedright
Element copy
\end{minipage} \\
\midrule\noalign{}
\endhead
\bottomrule\noalign{}
\endlastfoot
\texttt{from\_vec} & performed by caller; ownership moves & no \\
\texttt{zeros} & yes & zero initialization \\
\texttt{view} / \texttt{view\_mut} & no & no \\
\texttt{reshape\_view} & no & no \\
\texttt{transpose} & no & no \\
\texttt{slice\_axis} & no & no \\
\texttt{to\_contiguous} & yes & yes \\
\texttt{\seqsplit{OwnedTensor::clone}} & yes for nonempty payload & yes;
clones the owned vector \\
\texttt{\seqsplit{TensorView::clone}} & no & no; retains the storage
borrow \\
\end{longtable}
}

The distinction is shown in
\href{https://github.com/hermonai/inside-the-llm-engine/blob/astra-visual-rewrite/diagrams/tensor/view-vs-copy.txt}{view
versus copy}. Metadata allocation is intentionally separated from
element allocation in this table because model payload dominates later
memory costs.

The derived \texttt{Clone} implementation deserves explicit attention.
Cloning an owner duplicates its vectors, including the element payload;
it does not create a cheap shared handle. Cloning a view duplicates
metadata and the immutable reference, so both views still read the same
allocation. The operation's type determines which contract applies. A
review that searches only for calls named \texttt{copy} would miss the
owner's clone cost.

\EngineFigure{tensor-copy}{Materialization reads a strided logical sequence and writes a new canonical allocation.}{fig:tensor-copy}

For the transposed six-value view, \texttt{to\_contiguous} reads source
offsets \texttt{{[}0,3,1,4,2,5{]}} and writes destination offsets
\texttt{{[}0,1,2,3,4,5{]}}. The new physical buffer is
\texttt{{[}1,4,2,5,3,6{]}}. Its logical values equal the transposed
source even though its physical order and owner differ. Changing the
copy's \texttt{{[}0,1{]}} from 4 to 40 must leave the original matrix's
\texttt{{[}1,0{]}} equal to 4.

For \(N\) F32 elements, the analytical payload movement of this explicit
copy is

\begin{equation}
Q_{\mathrm{copy,payload}}=N\times4+N\times4=8N\quad[\mathrm{bytes}].
\end{equation}

Here the two terms count logical element reads and writes. For our
example, \(N=6\) gives 48 payload bytes. This is not a measurement of
DRAM traffic: cache lines, write allocation, metadata, allocator
bookkeeping and validation are outside the count. The function also
allocates coordinate vectors while enumerating logical indices. Its
clarity-first implementation is not a promise of a production memcpy
cost.
\href{https://github.com/hermonai/inside-the-llm-engine/blob/astra-visual-rewrite/figures/generated/tensor-copy.txt}{Text
sequence}.

\begin{quote}
\textbf{FIRST PRINCIPLE} Allocation and copying must be visible
operations. A convenient-looking metadata transform must not secretly
materialize a model-sized payload.
\end{quote}

\section{Reshape changes grouping}\label{reshape-changes-grouping}

A no-copy reshape changes how one contiguous sequence is grouped into
axes. The six physical values \texttt{1,2,3,4,5,6} can be viewed as
\texttt{{[}2,3{]}}, \texttt{{[}3,2{]}}, \texttt{{[}6{]}}, or
\texttt{{[}1,6{]}}. Logical row-major iteration remains physical offsets
\texttt{0,1,2,3,4,5}.

V1's \texttt{reshape\_view} requires two proofs:

\begin{enumerate}
\def\labelenumi{\arabic{enumi}.}
\tightlist
\item
  The source has exact canonical row-major strides.
\item
  The requested checked element count equals the current count.
\end{enumerate}

It then derives canonical strides for the new shape and borrows the same
storage. Shape \texttt{{[}4,2{]}} fails because it asks six values to
become eight. A non-contiguous transpose fails with
\texttt{NonContiguous}; v1 does not copy as a fallback. The canonical
\href{https://github.com/hermonai/inside-the-llm-engine/blob/astra-visual-rewrite/diagrams/tensor/reshape-view.txt}{reshape
view} makes the shared owner and changed grouping explicit.

\EngineFigure{tensor-reshape}{Reshape and transpose both produce three-by-two views but disagree at the same coordinate.}{fig:tensor-reshape}

Let \(\mathbf{R}\) be the \texttt{{[}3,2{]}} reshape of our original
owner. Its strides are \texttt{{[}2,1{]}}, and \(R_{0,1}=2\). The
transpose has shape \texttt{{[}3,2{]}} too, but its strides are
\texttt{{[}1,3{]}}, so \((A^{\mathsf T})_{0,1}=4\). Both results are
valid aliases. Checking only the shape would not reveal their different
mathematical values. An operator contract must therefore state whether
it honors general strides, requires canonical input, or performs an
explicit conversion.

PyTorch's \texttt{view} supports a broader compatibility condition based
on adjacent stride relationships, while \texttt{reshape} may choose a
view or copy. That is useful framework behavior. It is too implicit for
our first substrate, where the method name itself should reveal
allocation.

\section{Transpose changes logical
movement}\label{transpose-changes-logical-movement}

Transpose is not reshape. For the original \texttt{{[}2,3{]}} matrix:

\begin{equation}A=\begin{pmatrix}1&2&3\\4&5&6\end{pmatrix},\end{equation}

the rank-2 transpose is logically:

\begin{equation}A^\top=\begin{pmatrix}1&4\\2&5\\3&6\end{pmatrix}.\end{equation}

No bytes move. Shape \texttt{{[}2,3{]}} becomes \texttt{{[}3,2{]}};
strides \texttt{{[}3,1{]}} become \texttt{{[}1,3{]}}. Logical row-major
traversal now visits physical offsets \texttt{{[}0,3,1,4,2,5{]}}. The
canonical
\href{https://github.com/hermonai/inside-the-llm-engine/blob/astra-visual-rewrite/diagrams/tensor/transpose-view.txt}{transpose
view} shows both interpretations over one owner.

\EngineFigure{tensor-transpose}{Transpose swaps shape and stride axes while both views borrow the original six values.}{fig:tensor-transpose}

The coordinate identity is

\begin{equation}
(A^{\mathsf T})_{j,i}=A_{i,j},\qquad
\operatorname{offset}_{A^{\mathsf T}}(j,i)=j+3i
=\operatorname{offset}_A(i,j).
\end{equation}

For example, the transposed coordinate \texttt{{[}0,1{]}} and original
coordinate \texttt{{[}1,0{]}} both reach element 3, whose value is 4.
The Rust visual-contract test compares the returned references, not only
the numbers. Value equality could survive an accidental copy; reference
identity proves this particular view aliases the owner.
\href{https://github.com/hermonai/inside-the-llm-engine/blob/astra-visual-rewrite/figures/generated/tensor-transpose.txt}{Text
sequence}.

\begin{quote}
\textbf{ENGINEERING FAILURE --- TRANSPOSE ASSUMED CONTIGUOUS} A kernel
receives the transposed shape but ignores its strides. It reads physical
offsets \texttt{{[}0,1,2,3,4,5{]}}, producing finite, plausible, wrong
output. Layout is part of the kernel's semantics, not an optional
optimization hint.
\end{quote}

\section{Slices introduce a base
offset}\label{slices-introduce-a-base-offset}

Take rows \texttt{1..3} from a canonical \texttt{{[}4,3{]}} tensor. The
new shape is \texttt{{[}2,3{]}} and the strides remain
\texttt{{[}3,1{]}}, but the first view element is three elements into
the owner. Its base offset is 3.

\texttt{\seqsplit{slice\_axis(axis,start,end)}} implements only bounded
half-open ranges. It does not implement Python slicing syntax, steps,
reverse views, new axes, or fancy indices. The new base is

\begin{equation}
b' = b + \text{start}\times s_{\text{axis}},
\end{equation}

computed with checked arithmetic. The constructor then validates the
complete new extent. A zero-length slice may begin at the end of an axis
and reach no storage.

The representation keeps the full borrowed storage plus
\texttt{base\_offset} rather than slicing the Rust reference itself.
That makes the address formula and diagnostics explicit for later
model-file views. A subslice representation could encode part of the
bound in Rust's slice length, but would hide the original
storage-relative offset we want to teach.

\EngineFigure{tensor-slice}{A column slice retains gaps in storage and needs a larger backing range than its logical element count.}{fig:tensor-slice}

Return to the six-value matrix and select columns \texttt{1..3} using
\texttt{\seqsplit{slice\_axis(1,1,3)}}. The result is
\texttt{{[}{[}2,3{]},{[}5,6{]}{]}}. Its shape is \texttt{{[}2,2{]}}, its
strides remain \texttt{{[}3,1{]}}, and its base is 1. The offset formula
becomes

\begin{equation}
o_S(i,j)=1+3i+j,\qquad 0\le i<2,\quad 0\le j<2.
\end{equation}

The four coordinates reach offsets \texttt{{[}1,2,4,5{]}}. The logical
payload is four F32 values, or 16 bytes, but a borrowed storage slice
addressed from its own index zero must have length at least 6. It must
include the preceding element and the gap, even though those elements
are never logically read by this view. This is why ``four values'' and
``six backing slots'' are compatible facts. The test passes a
five-element backing slice with the same metadata and requires
\texttt{\seqsplit{InvalidViewExtent}} before any value can be read.

\section{Mutable aliasing is a numerical
hazard}\label{mutable-aliasing-is-a-numerical-hazard}

Two overlapping immutable views are safe. Two overlapping writable views
can race or make one operation observe another's partial output. The
result may be nondeterministic without ever becoming NaN or crashing.
Ownership is therefore a numerical correctness concern.

V1 does not offer a constructor for arbitrary mutable strides. An owner
can produce one \texttt{TensorViewMut} over its complete canonical
storage. That operation borrows the owner exclusively. While the mutable
view is live, Rust prevents another shared or mutable view from
coexisting.

\begin{verbatim}
let mut tensor = OwnedTensor::zeros(vec![2, 3])?;
{
    let mut writable = tensor.view_mut();
    *writable.get_mut(&[1, 2])? = 7.0;
}
assert_eq!(*tensor.view().get(&[1, 2])?, 7.0);
\end{verbatim}

\EngineFigure{tensor-lifetime}{Shared reads end before exclusive mutation, and a later shared view observes the new value.}{fig:tensor-lifetime}

In the recurring fixture, a shared read sees 6, its last use ends, an
exclusive view writes 60, and a later shared view sees 60. The
successful transition is covered by a normal Rust test. Two
\texttt{compile\_fail} documentation tests verify the forbidden cases:
using a shared view after requesting overlapping mutable access, and
returning a view of a local owner. These failures are checked by the
compiler as part of \texttt{cargo\ test\ -\/-doc}; no intentionally
broken file is placed in a normal test target.
\href{https://github.com/hermonai/inside-the-llm-engine/blob/astra-visual-rewrite/figures/generated/tensor-lifetime.txt}{Text
lifetime trace}.

A future safe split operation could prove that two canonical ranges do
not overlap before returning two mutable borrows. Chapter 5 does not
need it. It also needs no unsafe code: \texttt{engine0} retains
\texttt{\seqsplit{\#![forbid(unsafe\_code)]}}.

\section{Physical addresses, alignment, and
locality}\label{physical-addresses-alignment-and-locality}

\texttt{Vec\textless{}f32\textgreater{}} stores initialized \texttt{f32}
elements contiguously. If an illustrative base address were
\texttt{0x1000}, consecutive elements would begin at \texttt{0x1000},
\texttt{0x1004}, \texttt{0x1008}, and so on. Real allocator addresses
vary; the point is the four-byte displacement, not those invented
addresses.

The allocation is aligned suitably for \texttt{f32}. That does not
promise alignment for a future SIMD instruction, GPU buffer, stable C
ABI, or direct I/O request. Part VII will make arena and native
alignment explicit.

Layout also affects locality. CPU caches fetch regions larger than one
scalar, so visiting neighboring physical elements often benefits from
already-fetched data. For a row-major \texttt{{[}N,N{]}} matrix,
row-then-column iteration visits offsets \texttt{0,1,2,...}.
Column-then-row iteration visits \texttt{0,N,2N,...}, then
\texttt{1,N+1,2N+1,...}. The mathematical checksum can be identical
while the memory access pattern differs.

This is a preview, not a cache-architecture chapter. Associativity,
cache levels, hardware prefetching, NUMA, and kernel tiling remain
ahead.

\section{Build Tensor Substrate v1}\label{build-tensor-substrate-v1}

The implementation lives in \texttt{engine0::tensor}. A module is
enough: a separate crate would add a dependency boundary before another
consumer exists.

\begin{verbatim}
pub struct OwnedTensor {
    shape: Vec<usize>,
    strides: Vec<usize>,
    data: Vec<f32>,
}
\end{verbatim}

The owner stores canonical strides although they are derivable.
Constructors validate that redundancy once; repeated indexing and debug
output can then use the complete layout without recomputing metadata.
Dynamic rank was chosen over const generics because later operations
naturally move among vectors, matrices, and higher-rank activations. Our
purpose is to inspect runtime tensor metadata, not encode every rank
into a different Rust type.

The v1 public surface fits on one conceptual page:

\begin{itemize}
\tightlist
\item
  construct with \texttt{from\_vec} or \texttt{zeros};
\item
  inspect dtype, rank, shape, strides, length, and physical slice;
\item
  create canonical immutable or exclusive mutable views;
\item
  create a validated general immutable view from parts;
\item
  read through \texttt{get} or \texttt{get2};
\item
  create rank-2 transpose, bounded slice, or canonical reshape views;
\item
  explicitly materialize with \texttt{to\_contiguous};
\item
  call checked count, byte-count, stride, and offset helpers.
\end{itemize}

There is no broadcasting, autograd, device method, negative stride,
fancy indexing, generic operator graph, quantized dtype, BLAS dispatch,
SIMD, or GPU path. A tensor describes data. It does not own execution.

\section{Migrate ENGINE-1 without changing
it}\label{migrate-engine-1-without-changing-it}

Before this chapter, \texttt{\seqsplit{TinyLanguageModel}} stored
embedding and projection as bare vectors and indexed both manually. They
are now immutable \texttt{OwnedTensor} values with shape
\texttt{{[}V,D{]}} and strides \texttt{{[}D,1{]}}. The scalar forward
loop remains visible:

\begin{equation}
z_i=b_i+\sum_{j=0}^{D-1}W_{i,j}h_j,
\qquad
\mathbf{W}\in\mathbb{R}^{V_{\mathrm{vocab}}\times D},
\quad \mathbf{h}\in\mathbb{R}^{D},
\quad \mathbf{z}\in\mathbb{R}^{V_{\mathrm{vocab}}}.
\end{equation}

The model owns \(W\) with shape \texttt{{[}V,D{]}}, the request owns
\(h\) with shape \texttt{{[}D{]}}, and the forward result owns logits
\(z\) with shape \texttt{{[}V{]}}. Storage and accumulation are
\texttt{f32} in this tiny implementation. Each \texttt{W{[}i,j{]}} and
embedding element is now reached through checked \texttt{get2}; no
matrix multiplication operator has been smuggled into \texttt{Tensor}.

Bias, hidden activation, and \texttt{Logits} remain specialized
one-dimensional vectors. Wrapping every vector would increase
abstraction without eliminating an ambiguous multi-axis indexing
assumption. Chapter 6 can revisit activation representations when
operators need common tensor inputs.

This refactor is our first internal substitution gate. The
representation changed; model semantics did not. The full hand oracle
still gives \texttt{\seqsplit{[-0.7,0.1,0.4,2.2]}} for token
\texttt{like}, and greedy generation still emits \texttt{Rust} followed
by EOS.

\section{Prove the memory model
independently}\label{prove-the-memory-model-independently}

The plain-Python oracle does not import Rust code. It derives
\texttt{{[}12,4,1{]}} for shape \texttt{{[}2,3,4{]}}, maps
\texttt{{[}1,2,3{]}} to 23, enumerates the \texttt{{[}2,3{]}} matrix at
physical offsets \texttt{{[}0,1,2,3,4,5{]}}, and proves two different
\texttt{{[}3,2{]}} interpretations:

\begin{itemize}
\tightlist
\item
  reshape offsets: \texttt{{[}0,1,2,3,4,5{]}};
\item
  transpose offsets: \texttt{{[}0,3,1,4,2,5{]}}.
\end{itemize}

The transposed logical copy is exactly \texttt{{[}A,D,B,E,C,F{]}}. These
are integer and symbolic results, so no floating-point tolerance is
needed.

Run it from the repository root:

\begin{verbatim}
python3 code/reference/python/chapter05_tensor_oracle.py
\end{verbatim}

The Rust suite adds deterministic generated shapes of ranks one through
four. Every canonical logical index must map to its flat row-major
index, and the largest index must map to \texttt{element\_count\ -\ 1}.
Boundary tests exercise shape, byte, stride, offset, and extent overflow
without attempting giant allocations.

\section{Performance lab: same sum, different
walk}\label{performance-lab-same-sum-different-walk}

The Chapter 5 harness allocates one canonical \texttt{{[}2048,2048{]}}
tensor---4,194,304 \texttt{f32} values, or 16 MiB---and accumulates
every value as \texttt{f64}. It alternates row-major and column-wise
traversal across seven warmed release repetitions. Both orders produce
the exact checksum \texttt{524280621.0}.

On the recorded Apple M1 system, medians were:

{\def\LTcaptype{none} % do not increment counter
\begin{longtable}[]{@{}lr@{}}
\toprule\noalign{}
Order & Median \\
\midrule\noalign{}
\endhead
\bottomrule\noalign{}
\endlastfoot
Physical row-major & 4,163,875 ns \\
Column-wise logical & 13,692,583 ns \\
\end{longtable}
}

The durable
\href{https://github.com/hermonai/inside-the-llm-engine/blob/astra-visual-rewrite/research/benchmarks/chapter-05-traversal-order.md}{benchmark
record} contains the commit, command, hardware, software, shape,
repetitions, checksum, and raw result. It supports one narrow
observation: on that run, different access order had different cost. It
does not establish a universal ratio, isolate a cache level, or measure
matrix multiplication or inference.

\section{Follow the element, byte, and
owner}\label{follow-the-element-byte-and-owner}

One logical element now has a complete route:

A logical coordinate is validated against the shape, translated by strides and base offset, then read from the backing slice.

See
\href{https://github.com/hermonai/inside-the-llm-engine/blob/astra-visual-rewrite/diagrams/tensor/follow-the-element.txt}{Follow
the Element}. The route is where mathematical coordinates meet
memory-safe code.

The byte journey begins at an owned
\texttt{Vec\textless{}f32\textgreater{}}. Views, slices, reshapes, and
transposes change metadata while borrowing the same payload. Only
\texttt{to\_contiguous} creates a new allocation and moves logical
values into canonical order.
\href{https://github.com/hermonai/inside-the-llm-engine/blob/astra-visual-rewrite/diagrams/tensor/follow-the-byte.txt}{Follow
the Byte} records that fork.

The ownership journey begins with \texttt{OwnedTensor}, fans out into
immutable borrowed interpretations, and either ends those views before
the owner or creates an explicit independent copy.
\href{https://github.com/hermonai/inside-the-llm-engine/blob/astra-visual-rewrite/diagrams/tensor/follow-the-owner.txt}{Follow
the Owner} shows the lifetime boundary.

\section{A complete worked view}\label{a-complete-worked-view}

Let us run one allocation through every v1 metadata operation. Begin
with twelve values in canonical shape \texttt{{[}4,3{]}}:

\begin{equation}
X=
\begin{bmatrix}
0&1&2\\
3&4&5\\
6&7&8\\
9&10&11
\end{bmatrix}.
\end{equation}

\texttt{OwnedTensor::from\_vec({[}4,3{]},\ data)} calculates 12
elements, checks the data length, derives strides \texttt{{[}3,1{]}},
and takes ownership of the caller's vector. The separate
\texttt{\seqsplit{checked\_byte\_count}} helper establishes the 48-byte
payload; \texttt{from\_vec} does not call that helper. Its input is
already an initialized Rust vector, not an unchecked external byte
range. The logical coordinate \texttt{{[}2,1{]}} passes bounds \(2<4\)
and \(1<3\), then maps to

\begin{equation}
0 + 2\times3 + 1\times1 = 7.
\end{equation}

Storage element 7 contains value 7. No pointer arithmetic is exposed to
the caller.

Now take rows \texttt{1..3}. The slice has shape \texttt{{[}2,3{]}},
preserves strides \texttt{{[}3,1{]}}, and moves the base to

\begin{equation}
b'=0+1\times3=3.
\end{equation}

Its logical \texttt{{[}0,0{]}} reaches owner offset 3;
\texttt{{[}1,2{]}} reaches \(3+1\times3+2\times1=8\). The largest
reachable offset is 8, so a storage length of at least 9 is sufficient
for this view even though its owner contains 12 values. The view remains
canonical under v1's stride definition: shape \texttt{{[}2,3{]}} expects
\texttt{{[}3,1{]}}. It occupies one dense subregion from offsets 3
through 8.

Transpose that slice. Shape becomes \texttt{{[}3,2{]}}, strides become
\texttt{{[}1,3{]}}, and base remains 3. Logical \texttt{{[}2,1{]}} still
reaches offset

\begin{equation}
3+2\times1+1\times3=8,
\end{equation}

but logical iteration visits offsets \texttt{{[}3,6,4,7,5,8{]}}. This is
valid and non-contiguous. \texttt{\seqsplit{reshape\_view([6])}} refuses
it because v1 cannot reinterpret that logical walk as canonical physical
order without moving values.

Finally, \texttt{to\_contiguous} enumerates those six logical
coordinates, reads \texttt{{[}3,6,4,7,5,8{]}}, and creates a new
\texttt{{[}3,2{]}} owner whose physical storage holds exactly that
sequence with strides \texttt{{[}2,1{]}}. Source and destination have
equal logical values but different ownership and physical order.
Mutating the new owner cannot affect \(X\).

This worked path separates four statements that are easy to collapse:

\begin{itemize}
\tightlist
\item
  the original owner is contiguous;
\item
  the row slice is also contiguous but begins at a nonzero base;
\item
  the transpose is a valid non-contiguous alias;
\item
  the materialized result is a separate contiguous owner.
\end{itemize}

No single \texttt{shape} or \texttt{element\_count} field can prove all
four.

\section{Error taxonomy is part of the
API}\label{error-taxonomy-is-part-of-the-api}

A checked tensor API needs errors precise enough to identify which
invariant failed. V1 keeps the list small while preserving that
distinction.

{\def\LTcaptype{none} % do not increment counter
\begin{longtable}[]{@{}
  >{\raggedright\arraybackslash}p{(\linewidth - 4\tabcolsep) * \real{0.3333}}
  >{\raggedright\arraybackslash}p{(\linewidth - 4\tabcolsep) * \real{0.3333}}
  >{\raggedright\arraybackslash}p{(\linewidth - 4\tabcolsep) * \real{0.3333}}@{}}
\toprule\noalign{}
\begin{minipage}[b]{\linewidth}\raggedright
Error
\end{minipage} & \begin{minipage}[b]{\linewidth}\raggedright
Rejected condition
\end{minipage} & \begin{minipage}[b]{\linewidth}\raggedright
Why it is separate
\end{minipage} \\
\midrule\noalign{}
\endhead
\bottomrule\noalign{}
\endlastfoot
\texttt{ShapeOverflow} & element or canonical-stride product does not
fit & prevents undersized or unrepresentable owners \\
\texttt{\seqsplit{ByteCountOverflow}} & elements times dtype width does
not fit & element count can fit while bytes do not \\
\texttt{\seqsplit{StorageLengthMismatch}} & owner shape and supplied
vector length disagree & owner construction requires exact storage \\
\texttt{\seqsplit{ShapeStrideRankMismatch}} & view metadata vectors have
different lengths & no axis-to-stride mapping exists \\
\texttt{RankMismatch} & an index has too few or too many axes & logical
coordinate is malformed \\
\texttt{\seqsplit{IndexOutOfBounds}} & one coordinate is outside its
dimension & logical bounds failed before addressing \\
\texttt{OffsetOverflow} & extent or index arithmetic cannot be
represented & prevents wrapped physical addresses \\
\texttt{\seqsplit{InvalidViewExtent}} & reachable offsets exceed
borrowed storage & arbitrary strides need more than count validation \\
\texttt{NonContiguous} & reshape source lacks canonical strides &
no-copy contract cannot be met \\
\texttt{\seqsplit{ReshapeElementMismatch}} & requested count changes &
reshape cannot create or discard values \\
\texttt{ExpectedMatrix} & transpose source is not rank 2 & v1 implements
only the promised operation \\
\texttt{InvalidSlice} & axis or half-open range is malformed & slicing
stays bounded and explicit \\
\end{longtable}
}

The display text helps a human; callers can match the typed variant. No
checked operation converts a malformed request into a default tensor or
clamps an index. Silent recovery would turn structural corruption into
plausible numbers.

There are still internal invariants. An \texttt{OwnedTensor} has already
proven that its strides are canonical and its storage count is exact.
Methods may rely on that fact, but public inputs pass through checked
construction. This boundary lets code remain readable without repeating
an entire model-file parser's threat model at every loop iteration.

\section{Why the type system stops
here}\label{why-the-type-system-stops-here}

It is tempting to encode every fact in Rust types:

\begin{verbatim}
Tensor<Element, Device, Layout, Rank, Shape, Allocator, Mutability, ...>
\end{verbatim}

Such a design can be valuable in a specialized library. It would be
harmful at this point in the book. Transformer intermediates change rank
and shape at runtime. Model metadata eventually arrives from files. If
the reader spends the chapter decoding const-generic machinery, the
physical memory model has again disappeared behind an abstraction.

Dynamic \texttt{Vec\textless{}usize\textgreater{}} metadata makes each
runtime check visible. \texttt{F32} makes one dtype concrete. An
immutable general view demonstrates real striding. A narrow mutable view
demonstrates exclusivity. This is enough structure to prevent the
current errors and enough simplicity to audit on one page.

The substrate is also intentionally not an operator object. Adding
methods such as \texttt{\seqsplit{tensor.matmul(other)}} would couple
data representation to arithmetic, temporary allocation, provider
choice, and error policy before those contracts have been derived.
Chapter 6 will accept tensor views as inputs to separately named
kernels. Later a planner may choose a scalar CPU, SIMD, Metal, or CUDA
implementation while the input tensor remains data plus metadata.

\section{Correctness matrix}\label{correctness-matrix}

The test suite does more than exercise happy-path indexing. Each group
protects a later inference-engine boundary.

\textbf{Construction tests} cover rank 0 through rank 4, exact shape
retrieval, canonical strides, dtype width, storage mismatch, zero
dimensions, element overflow, byte overflow, and canonical-stride
overflow. The overflow cases call helpers directly near
\texttt{usize::MAX}; they never attempt dangerous allocations.

\textbf{Indexing tests} cover first, middle, and final elements, the
hand-computed \([1,2,3]\mapsto23\) case, wrong-rank coordinates,
per-axis bounds, base offsets, extent overflow, and strides that require
more storage than logical count. A generated loop enumerates every
coordinate for small ranks one through four and requires the canonical
physical offset sequence exactly.

\textbf{View tests} prove rank-2 transpose shape, strides, and values;
compatible reshape shapes; rejection of incompatible and non-contiguous
reshape; bounded axis slicing; zero-length slices; zero-stride immutable
aliasing; and explicit logical-order materialization. Reference identity
proves a view reaches the owner's element rather than a copied value.

\textbf{Ownership tests} mutate one complete canonical owner through
\texttt{TensorViewMut}, end the exclusive borrow, and read the new value
through a later immutable view. A separately materialized owner is
mutated and the source is proven unchanged. Compile-time borrow failures
belong in Lab 21 rather than normal test sources, because committed test
targets must compile.

\textbf{Regression tests} retain every Part I contract: tokenizer
boundaries, strict UTF-8 streaming, full numerical logits, sampling
distributions, seeded state, token feedback, cancellation, EOS, budgets,
failure paths, and exactly one terminal outcome. The tensor refactor is
invalid if any of those observable semantics changes.

Together these are stronger evidence than testing one printed matrix.
Shape, layout, storage, and lifetime are orthogonal enough that each
deserves its own failure injection.

\section{Review from four engineering
perspectives}\label{review-from-four-engineering-perspectives}

A systems programmer should now be able to locate every allocation and
copy. The model owns parameter vectors. Views borrow them. Metadata
transforms do not move elements. \texttt{to\_contiguous} does. Size,
extent, and offset arithmetic are checked at public boundaries.

An ML engineer should recognize standard rank, shape, stride, transpose,
and view concepts while seeing the declared differences from NumPy and
PyTorch. Model equations retain their shapes. No layout detail changes
the mathematical meaning of \(z=Wh+b\).

A Rust engineer should find no dangling reference escape and no safe
path to arbitrary overlapping mutable strided views. Shared borrows can
coexist; exclusive mutation cannot overlap them. There is no
\texttt{unsafe}, and malformed public indexing returns typed errors.

A beginner should be able to answer the concrete question ``where is
\texttt{{[}1,2{]}}?'': check two bounds, multiply each coordinate by its
element stride, add the base, then read that physical element. The term
\emph{tensor} now expands into inspectable data instead of shrinking
into framework vocabulary.

\section{Inside Hermon}\label{inside-hermon-1}

The following findings were re-verified on 2026-09-06 at Hermon commit
\texttt{5e908d7} and its llama.cpp submodule \texttt{389ff61d}; they are
not claims about an uninspected future revision.

\begin{quote}
\textbf{INSIDE HERMON --- LIBRARY CONTRACT} \texttt{hermon-llamacpp}
contains one of the project's unsafe FFI boundaries;
\texttt{hermon-kernels} also contains native FFI code. The bridge's safe
\texttt{\seqsplit{Model::tensor\_info}} copies a tensor's GGML
dimensions, rank, and type across the FFI boundary.
\texttt{tensor\_row\_f32} borrows the model handle for the call and
returns one newly owned decoded row. The shim validates tensor
existence, shape, host residency, layout, conversion support, and caller
output length.
\end{quote}

The loaded model owns packed tensor lifetime; Rust callers do not
receive a raw borrow that can outlive it. Unsafe handle access is
localized and each wrapper block states its invariant.
\texttt{TensorSession} similarly owns reusable mutable GGML scratch
state, is movable between threads, but is deliberately not
\texttt{Sync}; entry points require exclusive \texttt{\&mut} access.

\begin{quote}
\textbf{INSIDE HERMON --- LIBRARY} \texttt{hermon-gguf} exposes names,
dimensions, per-tensor serialized type, offsets, and checked byte
lengths. It can open an exactly bounded reader or read a checked
subrange without materializing the whole tensor. That proves parser
behavior, not that the default inference path uses Hermon-native
tensors.
\end{quote}

\begin{quote}
\textbf{INSIDE HERMON --- PREVIEW} The gated paged GGUF runtime combines
validated metadata with the llama.cpp packed tensor bridge. It keeps a
\texttt{TensorSession} behind a mutex for exclusive graph/scratch use.
The default production batched path remains llama.cpp managed, so
PREVIEW tensor plumbing must not be flattened into a CURRENT end-to-end
claim.
\end{quote}

Pinned GGML reinforces the same lesson. Its \texttt{ggml\_tensor}
records type, dimensions \texttt{ne}, byte strides \texttt{nb},
storage/buffer information, and view/source metadata because operations
cannot assume contiguity. GGML fixes a maximum rank and has packed dtype
concerns that v1 excludes. The shared principle is not API identity:
kernels need shape, layout, type, and lifetime facts together.

\EngineFigure{tensor-production}{Educational element strides and production byte-stride metadata require an explicit representation boundary.}{fig:tensor-production}

This comparison preserves the useful commonality without implying
identical APIs. Multiplying an F32 element stride by four gives a byte
displacement, but the same shortcut cannot describe arbitrary packed
quantized blocks. Similarly, a returned F32 row is a caller-owned
conversion result, not a Rust slice into packed model storage. The
\href{https://github.com/hermonai/inside-the-llm-engine/blob/astra-visual-rewrite/research/astra/chapter05-regeneration.md}{dated
source review} records the exact wrapper and runtime gates. Source
availability establishes these contracts; it does not establish a new
hardware benchmark or real-model equivalence result.

\section{Run the same tensor through every
representation}\label{run-the-same-tensor-through-every-representation}

The chapter's visual fixture is a contract that can fail independently
of the renderer. From the repository root, run:

\begin{verbatim}
python3 scripts/check-tensor-visual-parity.py
cd code/mini-engine
cargo test -p engine0 --test tensor_visual_contract
cargo test -p engine0 --doc
\end{verbatim}

The parity script asks a Rust example to construct the real owner and
call \texttt{transpose}, \texttt{to\_contiguous}, \texttt{reshape\_view}
and \texttt{slice\_axis}. It compares their reported shapes, strides,
bases, logical values and offsets with independent Python row/column
enumeration and the shared figure fixture. A hand-edited caption cannot
make an incorrect trace pass. Pointer and mutation tests then verify
what numbers alone cannot: aliases remain aliases, and copied/cloned
owners remain independent.

\Needspace{12\baselineskip}

{\def\LTcaptype{none} % do not increment counter
\begin{longtable}[]{@{}
  >{\raggedright\arraybackslash}p{(\linewidth - 8\tabcolsep) * \real{0.2000}}
  >{\raggedright\arraybackslash}p{(\linewidth - 8\tabcolsep) * \real{0.2000}}
  >{\raggedright\arraybackslash}p{(\linewidth - 8\tabcolsep) * \real{0.2000}}
  >{\raggedright\arraybackslash}p{(\linewidth - 8\tabcolsep) * \real{0.2000}}
  >{\raggedright\arraybackslash}p{(\linewidth - 8\tabcolsep) * \real{0.2000}}@{}}
\toprule\noalign{}
\begin{minipage}[b]{\linewidth}\raggedright
Representation
\end{minipage} & \begin{minipage}[b]{\linewidth}\raggedright
Shape
\end{minipage} & \begin{minipage}[b]{\linewidth}\raggedright
Element strides
\end{minipage} & \begin{minipage}[b]{\linewidth}\raggedright
Logical values
\end{minipage} & \begin{minipage}[b]{\linewidth}\raggedright
Owner
\end{minipage} \\
\midrule\noalign{}
\endhead
\bottomrule\noalign{}
\endlastfoot
Original A & \texttt{{[}2,3{]}} & \texttt{{[}3,1{]}} &
\texttt{{[}1,2,3,4,5,6{]}} & original \\
Transpose & \texttt{{[}3,2{]}} & \texttt{{[}1,3{]}} &
\texttt{{[}1,4,2,5,3,6{]}} & borrows original \\
Reshape & \texttt{{[}3,2{]}} & \texttt{{[}2,1{]}} &
\texttt{{[}1,2,3,4,5,6{]}} & borrows original \\
Materialized transpose & \texttt{{[}3,2{]}} & \texttt{{[}2,1{]}} &
\texttt{{[}1,4,2,5,3,6{]}} & new independent owner \\
Column slice, base 1 & \texttt{{[}2,2{]}} & \texttt{{[}3,1{]}} &
\texttt{{[}2,3,5,6{]}} & borrows original \\
\end{longtable}
}

This table is the chapter's synthesis: what entered was one allocation;
what we built was several checked interpretations and one explicit copy.
The next operator can now receive a truthful description of every input
it reads.

\section{Common mistakes}\label{common-mistakes-4}

\subsection{``A tensor is a GPU
object''}\label{a-tensor-is-a-gpu-object}

A tensor is a data/layout abstraction. V1 is entirely CPU-local. A later
provider may place tensor storage on a GPU, but device placement is not
the definition.

\subsection{``Tensor means matrix''}\label{tensor-means-matrix}

A matrix is the rank-2 case. Scalars, vectors, and higher-rank arrays
are also tensors under this book's definition.

\subsection{``Rank tells me linear
independence''}\label{rank-tells-me-linear-independence}

Tensor rank counts axes. Matrix rank measures a different algebraic
property.

\subsection{``Shape defines layout''}\label{shape-defines-layout}

Shape defines logical dimensions. Strides and base offset map those
dimensions to storage. Equal shape does not imply equal storage order.

\subsection{``The element count validates a
view''}\label{the-element-count-validates-a-view}

Four strided elements can reach offset 101. Required extent, not count
alone, must fit storage.

\subsection{``Strides are always
bytes''}\label{strides-are-always-bytes}

Their unit is an API contract. V1 uses elements; NumPy and GGML expose
bytes. Mixing the two multiplies every address error by dtype width.

\subsection{``Reshape and transpose are the
same''}\label{reshape-and-transpose-are-the-same}

Reshape regroups canonical order. Transpose changes which axis step
moves where and often becomes non-contiguous.

\subsection{``Transpose always
copies''}\label{transpose-always-copies}

V1 transpose only swaps metadata. An explicit \texttt{to\_contiguous}
performs a copy.

\subsection{``Views own their data''}\label{views-own-their-data}

A \texttt{TensorView\textless{}\textquotesingle{}a\textgreater{}} owns
metadata and borrows storage. It cannot outlive the owner. Two views may
alias.

\subsection{``A tensor should perform its own
matmul''}\label{a-tensor-should-perform-its-own-matmul}

Tensor is data plus metadata. A matrix multiplication kernel is an
operation with layout, arithmetic, and execution contracts. Chapter 6
keeps them separate.

\section{Labs}\label{labs-1}

Labs 16--21 turn each contract into evidence:

\begin{itemize}
\tightlist
\item
  \href{https://github.com/hermonai/inside-the-llm-engine/blob/astra-visual-rewrite/labs/lab-16-offset-by-hand.md}{Lab
  16} derives offsets by hand.
\item
  \href{https://github.com/hermonai/inside-the-llm-engine/blob/astra-visual-rewrite/labs/lab-17-transpose-without-copy.md}{Lab
  17} proves transpose aliases.
\item
  \href{https://github.com/hermonai/inside-the-llm-engine/blob/astra-visual-rewrite/labs/lab-18-reshape-view.md}{Lab
  18} checks no-copy reshape gates.
\item
  \href{https://github.com/hermonai/inside-the-llm-engine/blob/astra-visual-rewrite/labs/lab-19-non-contiguous-copy.md}{Lab
  19} materializes logical order.
\item
  \href{https://github.com/hermonai/inside-the-llm-engine/blob/astra-visual-rewrite/labs/lab-20-break-shape-arithmetic.md}{Lab
  20} attacks size and extent arithmetic.
\item
  \href{https://github.com/hermonai/inside-the-llm-engine/blob/astra-visual-rewrite/labs/lab-21-mutation-and-aliasing.md}{Lab
  21} uses exclusive mutation.
\end{itemize}

Each moves through CHECK, BUILD, BREAK, or EXTEND without adding later
Transformer operators.

\section{Exercises}\label{exercises-4}

\begin{enumerate}
\def\labelenumi{\arabic{enumi}.}
\tightlist
\item
  Derive strides for \texttt{{[}5{]}}, \texttt{{[}2,5{]}},
  \texttt{{[}3,2,5{]}}, and \texttt{{[}1,3,1{]}}. State units.
\item
  For shape \texttt{{[}4,3{]}}, slice rows \texttt{1..3}. Derive shape,
  strides, base, and largest reachable offset.
\item
  Give two different stride vectors for shape \texttt{{[}2,2{]}} over
  six storage values. List the logical visit order of each.
\item
  Explain why a zero-stride immutable view is safe to read but cannot be
  the basis of an arbitrary overlapping mutable API.
\item
  Calculate the minimum storage length for shape \texttt{{[}2,2{]}},
  strides \texttt{{[}100,1{]}}, base 7. Identify every overflow check
  the implementation needs.
\item
  Compare \texttt{{[}3,2{]}} reshape and transpose views of
  \texttt{{[}A,B,C,D,E,F{]}}. Which logical index first reveals that
  they differ?
\item
  Explain why \texttt{from\_vec} moves an allocation while
  \texttt{to\_contiguous} copies elements, even though both return an
  owner.
\item
  Design a typed error for a kernel that accepts only canonical rank-2
  input. Which checks belong to the tensor, and which belong to the
  kernel?
\item
  Sketch a safe API that splits a canonical matrix into two
  non-overlapping mutable row ranges. State the proof required before
  returning borrows.
\item
  Predict how element-stride APIs must adapt at an FFI boundary exposing
  byte strides and a packed quantized dtype.
\end{enumerate}

\section{What we still have not
built}\label{what-we-still-have-not-built}

Tensor Substrate v1 can describe and safely access data. It cannot
multiply two matrices, normalize an activation, form Q/K/V, encode
position, apply causal attention, run a feed-forward network, or stack
decoder layers. It has no broadcasting, training graph, packed dtype,
model loader, SIMD kernel, or accelerator provider.

Those exclusions are the point. A small substrate lets us audit every
invariant before performance work raises the stakes. We can now
distinguish a data structure error from an operator error and a
mathematical result from its physical traversal cost.

\section{Summary}\label{summary-4}

A tensor is not a vector with a suggestive variable name. It is storage
plus shape, layout, dtype, ownership, and lifetime. Rank counts logical
axes. Shape defines their dimensions. Checked products establish element
and byte counts. Strides and a base offset map bounded logical indices
to physical elements.

Owned tensors are canonical row-major \texttt{f32} storage in v1.
Immutable views may carry general non-negative strides and alias their
owner. Reshape borrows only canonical layouts with equal element counts.
Transpose swaps rank-2 metadata. Slices adjust one dimension and base.
\texttt{to\_contiguous} is the visible allocation and copy boundary.
Rust lifetimes prevent dangling views; exclusive borrowing keeps
arbitrary mutable aliasing out of the API. Every external size and
offset operation is checked.

ENGINE-1 now stores its embedding and projection as tensors, yet
produces the same logits and tokens. Representation changed without
changing model semantics. That is the substitution discipline the rest
of the book will need.

Return to the five opening questions. For ENGINE-1's embedding, the
answers are now executable facts: shape \texttt{{[}V,D{]}}; dtype
\texttt{F32}; canonical row-major strides \texttt{{[}D,1{]}}; storage
owned immutably by \texttt{\seqsplit{TinyLanguageModel}}; and views
bounded by a borrow of that model. For a transposed temporary, the shape
and strides change, the owner does not, and the borrow cannot escape.
For a contiguous copy, the logical values remain equal while ownership
and physical ordering change.

That checklist scales. When later chapters introduce activations, packed
model weights, KV blocks, and device buffers, each new representation
must answer the same questions before a kernel touches it. A familiar
equation is not a waiver. Neither is a familiar framework name. Correct
inference begins by making the data contract complete enough that wrong
layout, wrong lifetime, or wrong size cannot masquerade as valid
arithmetic.

\section{Chapter 6 preview: matrix
multiplication}\label{chapter-6-preview-matrix-multiplication}

We can represent matrices correctly. We still need to compute

\begin{equation}
\mathbf{C}=\mathbf{A}\mathbf{B}
\end{equation}

without treating the equation as an execution plan. It does not say
where \(A\), \(B\), and \(C\) live, which dimensions agree, which index
changes fastest, how many floating-point operations occur, how many
bytes move, or whether reuse fits a cache.

Chapter 6---\emph{Matrix Multiplication: The Engine Room}---will derive
dot product, matrix-vector multiplication, matrix-matrix multiplication,
loop ordering, operation count, memory traffic, arithmetic intensity,
tiling, reference kernels, and blocked scalar kernels on top of Tensor
Substrate v1. It will be the first chapter that deliberately thinks like
a performance engineer.

\section{References}\label{references-2}

\begin{itemize}
\tightlist
\item
  NumPy developers.
  \href{https://numpy.org/doc/stable/reference/generated/numpy.ndarray.html}{\texttt{ndarray}}
  and
  \href{https://numpy.org/doc/stable/reference/generated/numpy.ndarray.strides.html}{\texttt{ndarray.strides}}.
\item
  PyTorch contributors.
  \href{https://docs.pytorch.org/docs/stable/tensor_view.html}{Tensor
  Views},
  \href{https://docs.pytorch.org/docs/stable/generated/torch.Tensor.view.html}{\texttt{Tensor.view}},
  and
  \href{https://docs.pytorch.org/docs/stable/tensor_attributes.html}{Tensor
  Attributes}.
\item
The Rust Project.
  \href{https://doc.rust-lang.org/stable/reference/types/slice.html}{Slice
  types} and
  \href{https://doc.rust-lang.org/stable/book/ch04-02-references-and-borrowing.html}{References
  and Borrowing}.
\item
The Rust Project.
  \href{https://doc.rust-lang.org/std/primitive.usize.html\#method.checked_mul}{\texttt{\seqsplit{usize::checked\_mul}}}.
\item
  ggml-org. Pinned
  \href{https://github.com/ggml-org/llama.cpp/blob/389ff61d77b5c71cec0cf92fe4e5d01ace80b797/ggml/include/ggml.h}{\texttt{ggml.h}}.
\item
  Hermon source at commit
  \href{https://github.com/hermonai/hermon/commit/5e908d75e4d5ee102f97f6a880ed4d0968d949df}{\texttt{5e908d7}},
  especially \texttt{hermon-llamacpp}, \texttt{hermon-gguf}, and the
  paged runtime boundary.
\end{itemize}


\clearpage
\section{Worked synthesis: separate coordinates from addresses}
\subsection{Transpose without moving a byte}
\textbf{Problem.} Storage is $(1,2,3,4,5,6)$, shape $(2,3)$ and strides $(3,1)$.
Give the transposed view's shape, strides and value at $(2,1)$.

\textbf{Solution.} Swap metadata axes: shape $(3,2)$, strides $(1,3)$.
The requested offset is $2\cdot1+1\cdot3=5$, so the value is 6.
The new view still borrows the original storage. Its logical row-major
sequence is $(1,4,2,5,3,6)$; that is the sequence a materialized transpose
would write into a new canonical allocation.

\subsection{Logical count is not backing extent}
\textbf{Problem.} A view has shape $(3,)$, stride $(4,)$ and base 1.
How many logical values exist, and what minimum backing length is needed?

\textbf{Solution.} There are three logical values, at offsets $1,5,9$.
The backing slice must contain at least ten elements. Checking only that the
logical count 3 fits a slice of length 3 would allow out-of-bounds access.
For nonnegative strides, validate the maximum reachable offset with checked
arithmetic; handle zero extents separately because they reach no element.

\subsection{Ownership is stronger than an index calculation}
\textbf{Problem.} Why is a valid zero-stride read view not automatically a
valid mutable view?

\textbf{Solution.} Multiple logical indices can designate the same element.
Bounds validity proves reachability, not non-aliasing. Two mutable references
to overlapping elements would violate the exclusivity contract. The teaching
engine therefore keeps general views immutable and gates mutation on an
exclusive canonical owner.

